\chapter{Khoảng tin cậy và mô hình suy luận trong thí nghiệm chính}\label{app:main-ci}
Mỗi nhóm tuyến có nhiều chuyến và mỗi chuyến có thể được lấy mẫu nhiều lần. Để tính khoảng tin cậy 95\%, các quan sát được gộp trong từng nhóm trước khi lấy mẫu lại 10.000 lần trên 20 nhóm tuyến của tập kiểm tra. Với mỗi phương pháp và kịch bản, mô hình đạt Hit100 tốt nhất và mô hình có MAE thấp nhất được chọn riêng trên tập lựa chọn. Do đó, hai kết quả trong cùng một hàng có thể đến từ hai bộ ước lượng khác nhau.
\section{Độ chính xác khi suy luận vị trí}
\par\begingroup\footnotesize\setstretch{1.04}\setlength{\parskip}{0pt}
\begin{longtable}{@{}P{3.4cm}P{.7cm}P{2.624cm}P{2.624cm}P{2.624cm}P{2.624cm}@{}}
\caption{CI95\% đầy đủ cho năm tác vụ định vị/tái dựng}\label{tab:privacy-ci}\\
\toprule
\textbf{Phương pháp} & \textbf{S} & \textbf{Hit100\%} & \textbf{CI95\% Hit} & \textbf{MAE m} & \textbf{CI95\% MAE} \\
\midrule\endfirsthead
\multicolumn{6}{l}{\small\textit{Bảng \thetable\ (tiếp)}}\\\toprule
\textbf{Phương pháp} & \textbf{S} & \textbf{Hit100\%} & \textbf{CI95\% Hit} & \textbf{MAE m} & \textbf{CI95\% MAE} \\
\midrule\endhead
\midrule\multicolumn{6}{r}{\small\textit{Còn tiếp ở trang sau}}\\\endfoot
\bottomrule\endlastfoot
Raw GPS & S1 & 100{,}000 & [100{,}000; 100{,}000] & 0{,}0 & [0{,}0; 0{,}0] \\
Raw GPS & S2 & 100{,}000 & [100{,}000; 100{,}000] & 0{,}0 & [0{,}0; 0{,}0] \\
Raw GPS & S3 & 100{,}000 & [100{,}000; 100{,}000] & 0{,}0 & [0{,}0; 0{,}0] \\
Raw GPS & S9 & 100{,}000 & [100{,}000; 100{,}000] & 0{,}0 & [0{,}0; 0{,}0] \\
Raw GPS & S10 & 100{,}000 & [100{,}000; 100{,}000] & 0{,}0 & [0{,}0; 0{,}0] \\
Geo-I/REM Z (K1) & S1 & 18{,}522 & [14{,}892; 22{,}204] & 338{,}5 & [319{,}0; 359{,}8] \\
Geo-I/REM Z (K1) & S2 & 37{,}500 & [29{,}167; 45{,}833] & 214{,}3 & [171{,}2; 261{,}2] \\
Geo-I/REM Z (K1) & S3 & 18{,}522 & [14{,}892; 22{,}204] & 338{,}5 & [319{,}0; 359{,}8] \\
Geo-I/REM Z (K1) & S9 & 5{,}833 & [2{,}500; 10{,}000] & 435{,}8 & [380{,}2; 498{,}0] \\
Geo-I/REM Z (K1) & S10 & 35{,}000 & [26{,}667; 43{,}333] & 246{,}8 & [201{,}5; 297{,}8] \\
TransProtect* & S1 & 82{,}984 & [79{,}973; 85{,}672] & 54{,}4 & [47{,}9; 62{,}0] \\
TransProtect* & S2 & 94{,}167 & [88{,}333; 98{,}333] & 28{,}7 & [20{,}4; 39{,}1] \\
TransProtect* & S3 & 82{,}984 & [79{,}973; 85{,}672] & 54{,}4 & [47{,}9; 62{,}0] \\
TransProtect* & S9 & 96{,}667 & [92{,}500; 100{,}000] & 17{,}8 & [8{,}5; 29{,}3] \\
TransProtect* & S10 & 90{,}833 & [84{,}167; 96{,}667] & 33{,}6 & [23{,}2; 46{,}6] \\
Semantic correlation* & S1 & 82{,}258 & [77{,}097; 86{,}532] & 41{,}9 & [38{,}1; 46{,}1] \\
Semantic correlation* & S2 & 100{,}000 & [100{,}000; 100{,}000] & 8{,}2 & [5{,}1; 13{,}0] \\
Semantic correlation* & S3 & 82{,}258 & [77{,}097; 86{,}532] & 41{,}9 & [38{,}1; 46{,}1] \\
Semantic correlation* & S9 & 100{,}000 & [100{,}000; 100{,}000] & 5{,}3 & [5{,}0; 5{,}8] \\
Semantic correlation* & S10 & 100{,}000 & [100{,}000; 100{,}000] & 7{,}3 & [5{,}4; 10{,}6] \\
Fake-query insertion* & S1 & 53{,}387 & [41{,}156; 65{,}215] & 1213{,}0 & [907{,}7; 1539{,}9] \\
Fake-query insertion* & S2 & 92{,}500 & [85{,}000; 99{,}167] & 239{,}1 & [44{,}8; 481{,}7] \\
Fake-query insertion* & S3 & 98{,}790 & [97{,}957; 99{,}462] & 13{,}3 & [9{,}8; 18{,}5] \\
Fake-query insertion* & S9 & 99{,}167 & [97{,}500; 100{,}000] & 42{,}2 & [4{,}8; 116{,}4] \\
Fake-query insertion* & S10 & 100{,}000 & [100{,}000; 100{,}000] & 6{,}4 & [5{,}2; 7{,}9] \\
\textbf{GeoI-BQ v1} & S1 & 0{,}188 & [0{,}000; 0{,}565] & 2104{,}6 & [1784{,}5; 2436{,}0] \\
\textbf{GeoI-BQ v1} & S2 & 0{,}000 & [0{,}000; 0{,}000] & 2124{,}2 & [1789{,}4; 2462{,}4] \\
\textbf{GeoI-BQ v1} & S3 & 0{,}188 & [0{,}000; 0{,}565] & 2104{,}6 & [1784{,}5; 2436{,}0] \\
\textbf{GeoI-BQ v1} & S9 & 0{,}000 & [0{,}000; 0{,}000] & 1986{,}4 & [1626{,}2; 2343{,}0] \\
\textbf{GeoI-BQ v1} & S10 & 0{,}000 & [0{,}000; 0{,}000] & 2124{,}2 & [1789{,}4; 2462{,}4] \\
\end{longtable}\endgroup\par

\section{Chất lượng truy hồi theo mục đích tìm kiếm}
\Needspace{210pt}
\par\begingroup\small\setstretch{1.04}\setlength{\parskip}{0pt}
\begin{longtable}{@{}P{3.2cm}P{3.985cm}P{3.985cm}P{3.985cm}@{}}
\caption{Recall GeoI-BQ theo purpose và chênh lệch so với TransProtect*}\label{tab:purpose-ci}\\
\toprule
\textbf{Purpose} & \textbf{BQ Recall (\%)} & \textbf{CI95\% tuyệt đối} & \textbf{Δ so TransProtect* (pp), CI95\%} \\
\midrule\endfirsthead
\multicolumn{4}{l}{\small\textit{Bảng \thetable\ (tiếp)}}\\\toprule
\textbf{Purpose} & \textbf{BQ Recall (\%)} & \textbf{CI95\% tuyệt đối} & \textbf{Δ so TransProtect* (pp), CI95\%} \\
\midrule\endhead
\midrule\multicolumn{4}{r}{\small\textit{Còn tiếp ở trang sau}}\\\endfoot
\bottomrule\endlastfoot
Trung bình & 99{,}595 & [99{,}224; 99{,}903] & -0{,}401; [-0{,}773; -0{,}093] \\
Gần nhất & 99{,}638 & [99{,}280; 99{,}930] & -0{,}362; [-0{,}720; -0{,}070] \\
Nhanh nhất & 99{,}641 & [99{,}285; 99{,}930] & -0{,}359; [-0{,}715; -0{,}070] \\
Bán kính & 99{,}388 & [98{,}676; 99{,}894] & -0{,}612; [-1{,}324; -0{,}106] \\
Ít vòng đường & 99{,}712 & [99{,}362; 99{,}946] & -0{,}273; [-0{,}622; -0{,}040] \\
\end{longtable}\endgroup\par

\par\begingroup\small\setstretch{1.04}\setlength{\parskip}{0pt}
\begin{longtable}{@{}P{3.5cm}P{2.4cm}P{4.628cm}P{4.628cm}@{}}
\caption{Chênh lệch Recall theo purpose: GeoI-BQ trừ đối chứng}\label{tab:purpose-deltas}\\
\toprule
\textbf{Đối chứng} & \textbf{Purpose} & \textbf{BQ trừ control (pp)} & \textbf{CI95\% (pp)} \\
\midrule\endfirsthead
\multicolumn{4}{l}{\small\textit{Bảng \thetable\ (tiếp)}}\\\toprule
\textbf{Đối chứng} & \textbf{Purpose} & \textbf{BQ trừ control (pp)} & \textbf{CI95\% (pp)} \\
\midrule\endhead
\midrule\multicolumn{4}{r}{\small\textit{Còn tiếp ở trang sau}}\\\endfoot
\bottomrule\endlastfoot
Raw GPS & Trung bình & -0{,}404 & [-0{,}775; -0{,}095] \\
Raw GPS & Gần nhất & -0{,}362 & [-0{,}720; -0{,}070] \\
Raw GPS & Nhanh nhất & -0{,}359 & [-0{,}715; -0{,}070] \\
Raw GPS & Bán kính & -0{,}612 & [-1{,}324; -0{,}106] \\
Raw GPS & Ít vòng đường & -0{,}282 & [-0{,}634; -0{,}047] \\
Geo-I/REM Z (K1) & Trung bình & -0{,}105 & [-0{,}542; 0{,}275] \\
Geo-I/REM Z (K1) & Gần nhất & -0{,}114 & [-0{,}524; 0{,}241] \\
Geo-I/REM Z (K1) & Nhanh nhất & -0{,}112 & [-0{,}519; 0{,}242] \\
Geo-I/REM Z (K1) & Bán kính & -0{,}237 & [-1{,}006; 0{,}338] \\
Geo-I/REM Z (K1) & Ít vòng đường & 0{,}044 & [-0{,}381; 0{,}389] \\
TransProtect* & Trung bình & -0{,}401 & [-0{,}773; -0{,}093] \\
TransProtect* & Gần nhất & -0{,}362 & [-0{,}720; -0{,}070] \\
TransProtect* & Nhanh nhất & -0{,}359 & [-0{,}715; -0{,}070] \\
TransProtect* & Bán kính & -0{,}612 & [-1{,}324; -0{,}106] \\
TransProtect* & Ít vòng đường & -0{,}273 & [-0{,}622; -0{,}040] \\
Semantic correlation* & Trung bình & -0{,}404 & [-0{,}775; -0{,}096] \\
Semantic correlation* & Gần nhất & -0{,}362 & [-0{,}720; -0{,}070] \\
Semantic correlation* & Nhanh nhất & -0{,}359 & [-0{,}715; -0{,}070] \\
Semantic correlation* & Bán kính & -0{,}612 & [-1{,}324; -0{,}106] \\
Semantic correlation* & Ít vòng đường & -0{,}285 & [-0{,}637; -0{,}050] \\
Fake-query insertion* & Trung bình & -0{,}404 & [-0{,}775; -0{,}096] \\
Fake-query insertion* & Gần nhất & -0{,}362 & [-0{,}720; -0{,}070] \\
Fake-query insertion* & Nhanh nhất & -0{,}359 & [-0{,}715; -0{,}070] \\
Fake-query insertion* & Bán kính & -0{,}612 & [-1{,}324; -0{,}106] \\
Fake-query insertion* & Ít vòng đường & -0{,}285 & [-0{,}635; -0{,}051] \\
\end{longtable}\endgroup\par

\section{Bộ ước lượng được chọn cho từng kịch bản}
Bảng sau ghi cách suy ra vị trí từ các tọa độ công khai. Chẳng hạn, trọng tâm là trung bình các tọa độ trong một lần gửi. Với các phép ngoại suy ở đầu hoặc cuối chuyến, cần xác định thêm cách nối các điểm qua thời gian và số quan sát dùng để khớp đường thẳng. Các bộ ước lượng có học chỉ dùng dữ liệu huấn luyện. Nếu phương pháp chỉ công bố Q, bộ ước lượng không được thấy Z.
\par\begingroup\footnotesize\setstretch{1.04}\setlength{\parskip}{0pt}
\begin{longtable}{@{}P{3.1cm}P{.7cm}P{5.678cm}P{5.678cm}@{}}
\caption{Bộ ước lượng vị trí được chọn theo từng phương pháp và từng kịch bản}\label{tab:main-selected}\\
\toprule
\textbf{Phương pháp} & \textbf{S} & \textbf{Theo Hit100} & \textbf{Theo MAE} \\
\midrule\endfirsthead
\multicolumn{4}{l}{\small\textit{Bảng \thetable\ (tiếp)}}\\\toprule
\textbf{Phương pháp} & \textbf{S} & \textbf{Theo Hit100} & \textbf{Theo MAE} \\
\midrule\endhead
\midrule\multicolumn{4}{r}{\small\textit{Còn tiếp ở trang sau}}\\\endfoot
\bottomrule\endlastfoot
Raw GPS & S1 & Trọng tâm & Trọng tâm \\
Raw GPS & S2 & Trọng tâm & Trọng tâm \\
Raw GPS & S3 & Trọng tâm & Trọng tâm \\
Raw GPS & S9 & Trọng tâm & Trọng tâm \\
Raw GPS & S10 & Trọng tâm & Trọng tâm \\
Geo-I/REM Z (K1) & S1 & Trọng tâm & Trọng tâm \\
Geo-I/REM Z (K1) & S2 & Trọng tâm & Trọng tâm \\
Geo-I/REM Z (K1) & S3 & Trọng tâm & Trọng tâm \\
Geo-I/REM Z (K1) & S9 & Trọng tâm & Trọng tâm \\
Geo-I/REM Z (K1) & S10 & Trọng tâm & Quỹ đạo thứ nhất sau ghép hình học, ngoại suy từ 3 mốc \\
TransProtect* & S1 & Trọng tâm & Trọng tâm \\
TransProtect* & S2 & Trọng tâm & Trọng tâm \\
TransProtect* & S3 & Trọng tâm & Trọng tâm \\
TransProtect* & S9 & Trọng tâm & Quỹ đạo thứ nhất sau ghép hình học, ngoại suy từ 2 mốc \\
TransProtect* & S10 & Quỹ đạo thứ nhất sau ghép hình học, ngoại suy từ 6 mốc & Quỹ đạo thứ nhất sau ghép hình học, ngoại suy từ 6 mốc \\
Semantic correlation* & S1 & Trọng tâm & Điểm có xác suất tiên nghiệm cao nhất \\
Semantic correlation* & S2 & Trọng tâm & Trung vị theo từng tọa độ \\
Semantic correlation* & S3 & Trọng tâm & Điểm có xác suất tiên nghiệm cao nhất \\
Semantic correlation* & S9 & Trọng tâm & Điểm thứ 5 theo thứ tự gửi, ngoại suy từ 2 mốc \\
Semantic correlation* & S10 & Trọng tâm & Điểm thứ 5 ở lần gửi tại biên quan sát \\
Fake-query insertion* & S1 & Điểm có xác suất tiên nghiệm cao nhất & Hồi quy Extra Trees (64 cây) \\
Fake-query insertion* & S2 & Viterbi theo chuyển động & Viterbi theo chuyển động \\
Fake-query insertion* & S3 & Viterbi theo chuyển động & Viterbi theo chuyển động \\
Fake-query insertion* & S9 & Viterbi theo chuyển động, ngoại suy từ 2 mốc & Viterbi theo chuyển động, ngoại suy từ 2 mốc \\
Fake-query insertion* & S10 & Viterbi theo chuyển động, ngoại suy từ 2 mốc & Viterbi theo chuyển động \\
\textbf{GeoI-BQ v1} & S1 & Trọng tâm & Trung vị vị trí trên tập huấn luyện \\
\textbf{GeoI-BQ v1} & S2 & Trọng tâm & Trung vị vị trí trên tập huấn luyện \\
\textbf{GeoI-BQ v1} & S3 & Trọng tâm & Trung vị vị trí trên tập huấn luyện \\
\textbf{GeoI-BQ v1} & S9 & Trọng tâm & Trung vị vị trí trên tập huấn luyện (nhóm ghép hình học) \\
\textbf{GeoI-BQ v1} & S10 & Trọng tâm & Trung vị vị trí trên tập huấn luyện (nhóm ghép hình học) \\
\end{longtable}\endgroup\par


\chapter{Chi tiết các tác vụ bổ sung}\label{app:extended}
\section{S5: dự đoán nhánh đường trước và trong lúc rẽ}
Độ chính xác cho biết tỷ lệ chọn đúng nhánh trong hai nhánh có thể đi. Điểm Brier và mất mát logarit cho biết xác suất gán cho lựa chọn đó có phù hợp với kết quả quan sát hay không. Lần đầu xe đi vào cạnh đường cần đoán phải nằm sau thời điểm dự đoán. Ngay cả khi có GPS thật, độ chính xác vẫn gần mức đoán ngẫu nhiên ở cả hai pha. Khi bản thân dữ liệu quan sát đã ít phân biệt được hai nhánh, riêng tác vụ này khó cho thấy mức bảo vệ tăng thêm.
\par\begingroup\small\setstretch{1.04}\setlength{\parskip}{0pt}
\begin{longtable}{@{}P{3cm}P{2.1cm}P{7.775cm} P{1cm}P{1cm}@{}}
\caption{Tất cả pha S5, xác suất dự đoán và CI}\label{tab:s5-full}\\
\toprule
\textbf{Phương pháp} & \textbf{Pha} & \textbf{Acc\% [CI95\%]} & \textbf{Brier} & \textbf{Log loss} \\
\midrule\endfirsthead
\multicolumn{5}{l}{\small\textit{Bảng \thetable\ (tiếp)}}\\\toprule
\textbf{Phương pháp} & \textbf{Pha} & \textbf{Acc\% [CI95\%]} & \textbf{Brier} & \textbf{Log loss} \\
\midrule\endhead
\midrule\multicolumn{5}{r}{\small\textit{Còn tiếp ở trang sau}}\\\endfoot
\bottomrule\endlastfoot
Raw GPS & Trước fork & 50{,}00 [50{,}00; 50{,}00] & 0{,}250 & 0{,}693 \\
Raw GPS & Cửa sổ rẽ & 50{,}00 [50{,}00; 50{,}00] & 0{,}250 & 0{,}693 \\
REM-Z K1 & Trước fork & 50{,}00 [50{,}00; 50{,}00] & 0{,}250 & 0{,}693 \\
REM-Z K1 & Cửa sổ rẽ & 50{,}00 [50{,}00; 50{,}00] & 0{,}250 & 0{,}693 \\
TransProtect* & Trước fork & 50{,}00 [50{,}00; 50{,}00] & 0{,}250 & 0{,}693 \\
TransProtect* & Cửa sổ rẽ & 50{,}00 [50{,}00; 50{,}00] & 0{,}250 & 0{,}693 \\
Semantic correlation* & Trước fork & 50{,}00 [50{,}00; 50{,}00] & 0{,}250 & 0{,}693 \\
Semantic correlation* & Cửa sổ rẽ & 50{,}00 [50{,}00; 50{,}00] & 0{,}252 & 0{,}697 \\
Fake-query insertion* & Trước fork & 44{,}17 [38{,}33; 50{,}00] & 0{,}315 & 0{,}853 \\
Fake-query insertion* & Cửa sổ rẽ & 50{,}00 [43{,}33; 57{,}50] & 0{,}373 & 5{,}372 \\
\textbf{GeoI-BQ v1} & Trước fork & 52{,}50 [48{,}33; 57{,}50] & 0{,}259 & 0{,}712 \\
\textbf{GeoI-BQ v1} & Cửa sổ rẽ & 50{,}83 [45{,}83; 55{,}83] & 0{,}259 & 0{,}714 \\
\end{longtable}\endgroup\par

\section{S6: dự đoán điểm đến ở ba thời điểm}
Ba thời điểm quan sát nằm trước ngã rẽ, trong cửa sổ rẽ và ở lần gửi đầu tiên tính từ cửa sổ rẽ trở đi. Trong cả ba trường hợp, xe chỉ đến đích sau khi người quan sát đã đưa ra dự đoán.
\par\begingroup\footnotesize\setstretch{1.04}\setlength{\parskip}{0pt}
\begin{longtable}{@{}P{3cm}P{2.1cm}P{3.747cm} P{1cm}P{1cm}P{3.747cm}@{}}
\caption{Tất cả pha S6, xác suất dự đoán và sai số đích}\label{tab:s6-full}\\
\toprule
\textbf{Phương pháp} & \textbf{Pha} & \textbf{Acc\% [CI95\%]} & \textbf{Brier} & \textbf{Log loss} & \textbf{MAE m [CI95\%]} \\
\midrule\endfirsthead
\multicolumn{6}{l}{\small\textit{Bảng \thetable\ (tiếp)}}\\\toprule
\textbf{Phương pháp} & \textbf{Pha} & \textbf{Acc\% [CI95\%]} & \textbf{Brier} & \textbf{Log loss} & \textbf{MAE m [CI95\%]} \\
\midrule\endhead
\midrule\multicolumn{6}{r}{\small\textit{Còn tiếp ở trang sau}}\\\endfoot
\bottomrule\endlastfoot
Raw GPS & Trước fork & 50{,}00 [50{,}00; 50{,}00] & 0{,}250 & 0{,}693 & 572{,}7 [497{,}3; 647{,}8] \\
Raw GPS & Cửa sổ rẽ & 50{,}00 [50{,}00; 50{,}00] & 0{,}250 & 0{,}693 & 572{,}7 [497{,}3; 647{,}8] \\
Raw GPS & Sau rẽ & 100{,}00 [100{,}00; 100{,}00] & 0{,}034 & 0{,}181 & 2{,}4 [2{,}4; 2{,}4] \\
REM-Z K1 & Trước fork & 50{,}00 [50{,}00; 50{,}00] & 0{,}250 & 0{,}693 & 572{,}7 [497{,}3; 647{,}8] \\
REM-Z K1 & Cửa sổ rẽ & 50{,}00 [50{,}00; 50{,}00] & 0{,}250 & 0{,}693 & 572{,}7 [497{,}3; 647{,}8] \\
REM-Z K1 & Sau rẽ & 50{,}00 [50{,}00; 50{,}00] & 0{,}250 & 0{,}693 & 572{,}7 [497{,}3; 647{,}8] \\
TransProtect* & Trước fork & 50{,}00 [50{,}00; 50{,}00] & 0{,}250 & 0{,}693 & 572{,}7 [497{,}3; 647{,}8] \\
TransProtect* & Cửa sổ rẽ & 50{,}00 [50{,}00; 50{,}00] & 0{,}250 & 0{,}693 & 572{,}7 [497{,}3; 647{,}8] \\
TransProtect* & Sau rẽ & 95{,}00 [90{,}83; 98{,}33] & 0{,}168 & 0{,}518 & 63{,}5 [17{,}7; 119{,}5] \\
Semantic correlation* & Trước fork & 50{,}00 [50{,}00; 50{,}00] & 0{,}250 & 0{,}693 & 572{,}7 [497{,}3; 647{,}8] \\
Semantic correlation* & Cửa sổ rẽ & 50{,}00 [50{,}00; 50{,}00] & 0{,}252 & 0{,}697 & 572{,}7 [497{,}3; 647{,}8] \\
Semantic correlation* & Sau rẽ & 95{,}00 [87{,}50; 100{,}00] & 0{,}128 & 0{,}436 & 63{,}9 [2{,}4; 156{,}2] \\
Fake-query insertion* & Trước fork & 44{,}17 [38{,}33; 50{,}00] & 0{,}315 & 0{,}853 & 649{,}4 [528{,}1; 784{,}3] \\
Fake-query insertion* & Cửa sổ rẽ & 50{,}00 [43{,}33; 57{,}50] & 0{,}373 & 5{,}372 & 554{,}0 [461{,}6; 643{,}4] \\
Fake-query insertion* & Sau rẽ & 99{,}17 [97{,}50; 100{,}00] & 0{,}119 & 0{,}416 & 11{,}3 [2{,}4; 29{,}1] \\
\textbf{GeoI-BQ v1} & Trước fork & 52{,}50 [48{,}33; 57{,}50] & 0{,}259 & 0{,}712 & 540{,}5 [452{,}8; 627{,}4] \\
\textbf{GeoI-BQ v1} & Cửa sổ rẽ & 50{,}83 [45{,}83; 55{,}83] & 0{,}259 & 0{,}714 & 558{,}5 [469{,}3; 646{,}7] \\
\textbf{GeoI-BQ v1} & Sau rẽ & 50{,}83 [44{,}17; 57{,}50] & 0{,}259 & 0{,}712 & 568{,}9 [458{,}8; 677{,}0] \\
\end{longtable}\endgroup\par

\section{S8: tác động của kênh quan sát thứ hai}
Chênh lệch được tính so với khi chỉ quan sát người dùng mục tiêu. Đơn vị là điểm phần trăm đối với Hit100 và mét đối với MAE. Khoảng tin cậy được tính bằng cách lấy mẫu lại 10 khối, mỗi khối gồm hai nhóm tuyến, để giữ quan hệ phụ thuộc khi ghép với một người đi tuyến khác. Nếu mọi lần đo đều trượt, khoảng [0;0] chỉ mô tả mẫu đã quan sát; xác suất định vị thành công trong thực tế vẫn có thể lớn hơn 0.
\par\begingroup\footnotesize\setstretch{1.04}\setlength{\parskip}{0pt}
\begin{longtable}{@{}P{2.8cm}P{1.5cm}P{2.574cm}P{2.574cm}P{2.574cm}P{2.574cm}@{}}
\caption{S8 V2: CI và chênh lệch ghép theo block}\label{tab:s8-full}\\
\toprule
\textbf{Phương pháp} & \textbf{View} & \textbf{Hit\% [CI]} & \textbf{MAE m [CI]} & \textbf{ΔHit pp [CI]} & \textbf{ΔMAE m [CI]} \\
\midrule\endfirsthead
\multicolumn{6}{l}{\small\textit{Bảng \thetable\ (tiếp)}}\\\toprule
\textbf{Phương pháp} & \textbf{View} & \textbf{Hit\% [CI]} & \textbf{MAE m [CI]} & \textbf{ΔHit pp [CI]} & \textbf{ΔMAE m [CI]} \\
\midrule\endhead
\midrule\multicolumn{6}{r}{\small\textit{Còn tiếp ở trang sau}}\\\endfoot
\bottomrule\endlastfoot
Raw GPS & Target & 100{,}00\newline [100{,}00; 100{,}00] & 0{,}0\newline [0{,}0; 0{,}0] & 0{,}00\newline [0{,}00; 0{,}00] & 0{,}0\newline [0{,}0; 0{,}0] \\
Raw GPS & + protected & 100{,}00\newline [100{,}00; 100{,}00] & 0{,}0\newline [0{,}0; 0{,}0] & 0{,}00\newline [0{,}00; 0{,}00] & 0{,}0\newline [0{,}0; 0{,}0] \\
Raw GPS & + raw & 100{,}00\newline [100{,}00; 100{,}00] & 0{,}0\newline [0{,}0; 0{,}0] & 0{,}00\newline [0{,}00; 0{,}00] & 0{,}0\newline [0{,}0; 0{,}0] \\
Raw GPS & + unrelated & 100{,}00\newline [100{,}00; 100{,}00] & 0{,}0\newline [0{,}0; 0{,}0] & 0{,}00\newline [0{,}00; 0{,}00] & 0{,}0\newline [0{,}0; 0{,}0] \\
REM-Z K1 & Target & 10{,}83\newline [5{,}00; 16{,}67] & 359{,}7\newline [297{,}6; 424{,}9] & 0{,}00\newline [0{,}00; 0{,}00] & 0{,}0\newline [0{,}0; 0{,}0] \\
REM-Z K1 & + protected & 6{,}67\newline [0{,}00; 15{,}00] & 300{,}8\newline [247{,}9; 357{,}9] & -4{,}17\newline [-11{,}67; 2{,}50] & -58{,}9\newline [-74{,}7; -43{,}8] \\
REM-Z K1 & + raw & 100{,}00\newline [100{,}00; 100{,}00] & 0{,}0\newline [0{,}0; 0{,}0] & 89{,}17\newline [83{,}33; 95{,}00] & -359{,}7\newline [-424{,}9; -297{,}6] \\
REM-Z K1 & + unrelated & 10{,}83\newline [5{,}00; 16{,}67] & 359{,}7\newline [297{,}6; 424{,}9] & 0{,}00\newline [0{,}00; 0{,}00] & 0{,}0\newline [0{,}0; 0{,}0] \\
TransProtect* & Target & 60{,}83\newline [51{,}67; 70{,}00] & 102{,}8\newline [84{,}8; 121{,}9] & 0{,}00\newline [0{,}00; 0{,}00] & 0{,}0\newline [0{,}0; 0{,}0] \\
TransProtect* & + protected & 64{,}17\newline [56{,}67; 71{,}67] & 96{,}8\newline [78{,}8; 116{,}2] & 3{,}33\newline [0{,}00; 6{,}67] & -5{,}9\newline [-8{,}5; -3{,}6] \\
TransProtect* & + raw & 98{,}33\newline [95{,}83; 100{,}00] & 0{,}0\newline [0{,}0; 0{,}0] & 37{,}50\newline [27{,}50; 47{,}50] & -102{,}8\newline [-121{,}9; -84{,}8] \\
TransProtect* & + unrelated & 60{,}83\newline [51{,}67; 70{,}00] & 102{,}8\newline [84{,}8; 121{,}9] & 0{,}00\newline [0{,}00; 0{,}00] & 0{,}0\newline [0{,}0; 0{,}0] \\
Semantic correlation* & Target & 45{,}00\newline [20{,}00; 70{,}00] & 101{,}4\newline [76{,}7; 132{,}1] & 0{,}00\newline [0{,}00; 0{,}00] & 0{,}0\newline [0{,}0; 0{,}0] \\
Semantic correlation* & + protected & 75{,}00\newline [60{,}00; 90{,}00] & 46{,}2\newline [29{,}3; 62{,}5] & 30{,}00\newline [10{,}00; 50{,}00] & -55{,}3\newline [-95{,}8; -23{,}9] \\
Semantic correlation* & + raw & 100{,}00\newline [100{,}00; 100{,}00] & 0{,}0\newline [0{,}0; 0{,}0] & 55{,}00\newline [30{,}00; 80{,}00] & -101{,}4\newline [-132{,}1; -76{,}7] \\
Semantic correlation* & + unrelated & 45{,}00\newline [20{,}00; 70{,}00] & 101{,}4\newline [76{,}7; 132{,}1] & 0{,}00\newline [0{,}00; 0{,}00] & 0{,}0\newline [0{,}0; 0{,}0] \\
Fake-query insertion* & Target & 16{,}67\newline [5{,}00; 31{,}67] & 1857{,}4\newline [1526{,}6; 2225{,}8] & 0{,}00\newline [0{,}00; 0{,}00] & 0{,}0\newline [0{,}0; 0{,}0] \\
Fake-query insertion* & + protected & 68{,}33\newline [56{,}67; 80{,}00] & 1857{,}4\newline [1526{,}6; 2225{,}8] & 51{,}67\newline [46{,}67; 56{,}67] & 0{,}0\newline [0{,}0; 0{,}0] \\
Fake-query insertion* & + raw & 100{,}00\newline [100{,}00; 100{,}00] & 0{,}0\newline [0{,}0; 0{,}0] & 83{,}33\newline [68{,}33; 95{,}00] & -1857{,}4\newline [-2225{,}8; -1526{,}6] \\
Fake-query insertion* & + unrelated & 16{,}67\newline [5{,}00; 31{,}67] & 1857{,}4\newline [1526{,}6; 2225{,}8] & 0{,}00\newline [0{,}00; 0{,}00] & 0{,}0\newline [0{,}0; 0{,}0] \\
\textbf{GeoI-BQ v1} & Target & 0{,}00\newline [0{,}00; 0{,}00] & 2246{,}7\newline [2036{,}7; 2468{,}2] & 0{,}00\newline [0{,}00; 0{,}00] & 0{,}0\newline [0{,}0; 0{,}0] \\
\textbf{GeoI-BQ v1} & + protected & 0{,}00\newline [0{,}00; 0{,}00] & 2226{,}8\newline [1950{,}9; 2507{,}3] & 0{,}00\newline [0{,}00; 0{,}00] & -19{,}9\newline [-161{,}8; 111{,}7] \\
\textbf{GeoI-BQ v1} & + raw & 100{,}00\newline [100{,}00; 100{,}00] & 0{,}0\newline [0{,}0; 0{,}0] & 100{,}00\newline [100{,}00; 100{,}00] & -2246{,}7\newline [-2468{,}2; -2036{,}7] \\
\textbf{GeoI-BQ v1} & + unrelated & 0{,}00\newline [0{,}00; 0{,}00] & 2246{,}7\newline [2036{,}7; 2468{,}2] & 0{,}00\newline [0{,}00; 0{,}00] & 0{,}0\newline [0{,}0; 0{,}0] \\
\end{longtable}\endgroup\par

\Needspace{210pt}
\par\begingroup\small\setstretch{1.04}\setlength{\parskip}{0pt}
\begin{longtable}{@{}P{3.2cm}P{2.919cm}P{2.919cm}P{2.919cm}P{2.919cm}@{}}
\caption{Số tọa độ attacker quan sát tại cutoff S8, tính cả hai kênh}\label{tab:s8-observations}\\
\toprule
\textbf{Phương pháp} & \textbf{Target} & \textbf{+ protected} & \textbf{+ raw} & \textbf{+ unrelated} \\
\midrule\endfirsthead
\multicolumn{5}{l}{\small\textit{Bảng \thetable\ (tiếp)}}\\\toprule
\textbf{Phương pháp} & \textbf{Target} & \textbf{+ protected} & \textbf{+ raw} & \textbf{+ unrelated} \\
\midrule\endhead
\midrule\multicolumn{5}{r}{\small\textit{Còn tiếp ở trang sau}}\\\endfoot
\bottomrule\endlastfoot
Raw GPS & 6{,}2 & 12{,}4 & 12{,}4 & 12{,}4 \\
REM-Z K1 & 6{,}2 & 12{,}4 & 12{,}4 & 12{,}4 \\
TransProtect* & 6{,}2 & 12{,}4 & 12{,}4 & 12{,}4 \\
Semantic correlation* & 31{,}0 & 62{,}0 & 37{,}2 & 62{,}0 \\
Fake-query insertion* & 36{,}8 & 73{,}5 & 43{,}0 & 73{,}5 \\
\textbf{GeoI-BQ v1} & 31{,}0 & 62{,}0 & 37{,}2 & 62{,}0 \\
\end{longtable}\endgroup\par

\section{Bộ ước lượng trong các tác vụ bổ sung của GeoI-BQ}
Ở S5 và S6, một cách phân biệt hai nhánh là đo khoảng cách từ các điểm Q tới hình dạng công khai của từng nhánh rẽ. Tham số nhiệt độ điều chỉnh mức chênh lệch xác suất giữa hai lựa chọn. Bảng ghi các bộ ước lượng được chọn theo từng tác vụ. Một bộ ước lượng khác ngoài các phương án đã thử vẫn có thể suy luận tốt hơn.
\Needspace{210pt}
\par\begingroup\small\setstretch{1.04}\setlength{\parskip}{0pt}
\begin{longtable}{@{}P{3cm}P{12.719cm}@{}}
\caption{Các bộ ước lượng được chọn khi đánh giá GeoI-BQ trên các tác vụ bổ sung}\label{tab:extra-selected}\\
\toprule
\textbf{Tác vụ} & \textbf{Bộ ước lượng được chọn} \\
\midrule\endfirsthead
\multicolumn{2}{l}{\small\textit{Bảng \thetable\ (tiếp)}}\\\toprule
\textbf{Tác vụ} & \textbf{Bộ ước lượng được chọn} \\
\midrule\endhead
\midrule\multicolumn{2}{r}{\small\textit{Còn tiếp ở trang sau}}\\\endfoot
\bottomrule\endlastfoot
S4, liên kết phiên & Extra Trees (96 cây, ngưỡng 0,5) \\
S5, lúc rẽ & Khoảng cách nhỏ nhất tới nhánh rẽ (nhiệt độ 0,02) \\
S6, sau rẽ & Khoảng cách nhỏ nhất tới nhánh rẽ (nhiệt độ 0,02) \\
S8, Chỉ mục tiêu & Hit100: Trọng tâm. MAE: Hồi quy 5 láng giềng gần nhất. \\
S8, Kênh đã bảo vệ & Hit100: Trọng tâm. MAE: Hồi quy 5 láng giềng gần nhất. \\
S8, Kênh GPS thật & Hit100 và MAE: Trọng tâm trên kênh người đi cùng \\
S8, Người khác tuyến & Hit100: Trọng tâm. MAE: 5 láng giềng gần nhất, chỉ dùng kênh mục tiêu. \\
\end{longtable}\endgroup\par


\chapter{Ảnh hưởng của cách chọn Q và độ sâu phản hồi}\label{app:support}
\section{Vai trò của từng thành phần khi chọn Q}
Để xem việc chọn Q theo thông tin vị trí đã bảo vệ đóng góp bao nhiêu, thí nghiệm ngày 14/10 lần lượt thay đổi cách chọn tọa độ và giữ lại các thành phần còn lại theo cùng giao thức. Các hàng trong bảng dưới là những cấu hình nội bộ. Phép đối chiếu với ba phương pháp từ tài liệu nghiên cứu nằm ở Chương~\ref{ch:experiments}.
\Needspace{210pt}
\par\begingroup\small\setstretch{1.04}\setlength{\parskip}{0pt}
\begin{longtable}{@{}P{4cm}P{2.119cm}P{2.119cm}P{2.119cm}P{2.119cm}P{2.119cm}@{}}
\caption{Ablation nội bộ ngày 14/10, tách khỏi đối chiếu paper}\label{tab:ablation}\\
\toprule
\textbf{Ablation14/10} & \textbf{Trung bình} & \textbf{Gần nhất} & \textbf{Nhanh nhất} & \textbf{Bán kính} & \textbf{Ít vòng đường} \\
\midrule\endfirsthead
\multicolumn{6}{l}{\small\textit{Bảng \thetable\ (tiếp)}}\\\toprule
\textbf{Ablation14/10} & \textbf{Trung bình} & \textbf{Gần nhất} & \textbf{Nhanh nhất} & \textbf{Bán kính} & \textbf{Ít vòng đường} \\
\midrule\endhead
\midrule\multicolumn{6}{r}{\small\textit{Còn tiếp ở trang sau}}\\\endfoot
\bottomrule\endlastfoot
Service-cover K5 & 100{,}000 & 100{,}000 & 100{,}000 & 100{,}000 & 100{,}000 \\
Initial-only & 98{,}377 & 98{,}984 & 98{,}987 & 96{,}358 & 99{,}181 \\
Local explore & 99{,}562 & 99{,}609 & 99{,}609 & 99{,}358 & 99{,}670 \\
Public-base $\Gamma=0$ & 99{,}711 & 99{,}730 & 99{,}731 & 99{,}642 & 99{,}740 \\
GeoI-BQ v1 & 99{,}595 & 99{,}638 & 99{,}641 & 99{,}388 & 99{,}712 \\
\end{longtable}\endgroup\par

Trên các tuyến đã xét, cách chỉ dùng thông tin công khai đạt Recall cao hơn cấu hình $\Gamma=1$. Kết quả này minh họa sự khác biệt giữa điểm dùng để chọn đoạn Q và chất lượng của danh sách địa điểm cuối cùng.
\section{Khi giảm số địa điểm trong mỗi phản hồi}
Dịch vụ trả tối đa 60 địa điểm mỗi loại tại mỗi tọa độ. Để đo ảnh hưởng của độ sâu phản hồi, thí nghiệm giữ nguyên tọa độ và thời điểm gửi, rồi chỉ lấy 30 địa điểm đầu. Khoảng tin cậy cho chênh lệch L60--L30 được tính theo cặp trên cùng nhóm tuyến.
\Needspace{210pt}
\par\begingroup\small\setstretch{1.04}\setlength{\parskip}{0pt}
\begin{longtable}{@{}P{3cm}P{4.052cm}P{4.052cm}P{4.052cm}@{}}
\caption{Truncation L30 so với L60 trên cùng Q/time}\label{tab:depth}\\
\toprule
\textbf{Purpose} & \textbf{L30 cùng Q (\%)} & \textbf{L60 trừ L30 (pp)} & \textbf{CI95\% (pp)} \\
\midrule\endfirsthead
\multicolumn{4}{l}{\small\textit{Bảng \thetable\ (tiếp)}}\\\toprule
\textbf{Purpose} & \textbf{L30 cùng Q (\%)} & \textbf{L60 trừ L30 (pp)} & \textbf{CI95\% (pp)} \\
\midrule\endhead
\midrule\multicolumn{4}{r}{\small\textit{Còn tiếp ở trang sau}}\\\endfoot
\bottomrule\endlastfoot
Trung bình & 94{,}170 & 5{,}425 & [3{,}486; 7{,}527] \\
Gần nhất & 96{,}105 & 3{,}534 & [2{,}371; 4{,}726] \\
Nhanh nhất & 96{,}101 & 3{,}540 & [2{,}375; 4{,}740] \\
Bán kính & 89{,}015 & 10{,}374 & [5{,}768; 15{,}706] \\
Ít vòng đường & 95{,}459 & 4{,}253 & [2{,}946; 5{,}520] \\
\end{longtable}\endgroup\par


\chapter{Phạm vi dữ liệu và cách đối chiếu kết quả}\label{app:repro}
\section{Nguồn của các phép so sánh}
Thí nghiệm chính đối chiếu sáu phương pháp trên năm tác vụ định vị, cho 30 cặp phương pháp và kịch bản. Mỗi phương pháp được chấm trên cùng nhóm tuyến và cửa sổ thời gian. Với GeoI-BQ, các điểm Q và thời điểm gửi đã được cố định trước khi đối chiếu; cùng một chuỗi được dùng để tính độ riêng tư, chất lượng tìm kiếm và chi phí.

Đánh giá bổ sung gồm 66 tác vụ, trong đó có phép kiểm tra nội dung truy vấn ở S7 và lần đánh giá đầu tiên ở S8. Riêng S8 có thêm 24 tác vụ để xét những cách kết hợp thông tin từ hai người dùng. Khi thiếu nhãn thực tế để chấm một mục tiêu, kết quả được ghi là N/A. Các thí nghiệm về độ sâu phản hồi và về cách chọn Q trả lời những câu hỏi riêng. Mỗi chênh lệch được tính trong phạm vi thí nghiệm có chung điều kiện so sánh.
\section{Số chuyến được đánh giá}
Cả sáu phương pháp đều hoàn tất 120 trên 120 lượt đánh giá trong tập kiểm tra. Con số này cho biết có bao nhiêu quan sát để tính các chỉ số. Mức bảo vệ còn phụ thuộc vào người quan sát thấy gì và họ có thể suy ra vị trí chính xác đến đâu.
\Needspace{210pt}
\par\begingroup\small\setstretch{1.04}\setlength{\parskip}{0pt}
\begin{longtable}{@{}P{4cm}P{3.719cm}P{3.719cm}P{3.719cm}@{}}
\caption{Coverage TEST và số lần thực thi thất bại của trial chính}\label{tab:coverage}\\
\toprule
\textbf{Phương pháp} & \textbf{TEST thử} & \textbf{Hoàn tất} & \textbf{Thất bại} \\
\midrule\endfirsthead
\multicolumn{4}{l}{\small\textit{Bảng \thetable\ (tiếp)}}\\\toprule
\textbf{Phương pháp} & \textbf{TEST thử} & \textbf{Hoàn tất} & \textbf{Thất bại} \\
\midrule\endhead
\midrule\multicolumn{4}{r}{\small\textit{Còn tiếp ở trang sau}}\\\endfoot
\bottomrule\endlastfoot
Raw GPS & 120 & 120 & 0 \\
Geo-I/REM Z (K1) & 120 & 120 & 0 \\
TransProtect* & 120 & 120 & 0 \\
Semantic correlation* & 120 & 120 & 0 \\
Fake-query insertion* & 120 & 120 & 0 \\
\textbf{GeoI-BQ v1} & 120 & 120 & 0 \\
\end{longtable}\endgroup\par

\section{Đọc hai kết quả đối chứng ở S8}
Khi người đi cùng công bố GPS thật, lấy trọng tâm tọa độ trên kênh đó cho Hit100 bằng 100\% trong điều kiện thí nghiệm. Ban đầu, tiêu chí kiểm tra yêu cầu mọi mô hình được chọn cũng phải đạt mức này. Kết quả của TransProtect* cho thấy điều kiện ấy quá mạnh. Một bộ ước lượng khác đồng hạng trên tập lựa chọn được dùng để đánh giá và đạt 98,33\% trên tập kiểm tra. Đối chiếu lại từng dự đoán xác nhận cả hai giá trị. Khi đánh giá rủi ro tiết lộ, phép đọc trực tiếp cho thấy thông tin mà kênh thứ hai thực sự cung cấp, kể cả khi mô hình được chọn chưa tận dụng hết thông tin đó.
